%=============================================================================!
% 18.9  WHAT THE TRAINING LOOKS LIKE                                          !
%=============================================================================!
\section{What the training looks like}
\label{sec:lo-training}

A run of molecular dynamics is read through its energies, and a training
run through its losses. The curve of the loss against the step, for the
training readings and for the validation readings, is the first thing to
look at when a model is trained, and most faults show in it.
\Cref{fig:lo-training} trains the polynomials of \cref{sec:lo-basis} on
the fifteen readings of argon's pair energy by lowering the mean squared
miss from zero weights, once healthily and then broken in four ways, and
\cref{tab:lo-faults} gathers the checks.

\textit{Healthy.} Degree 8 in the scaled distance, Adam at the rate
$10^{-4}$, 40\,000 steps: the training miss falls to
\SI{0.202}{\milli\electronvolt} and the validation miss with it, to
0.187, both at the noise of \SI{0.2}{\milli\electronvolt}
(\cref{fig:lo-training}a). Two curves that fall together and level at the
noise are the picture of a model that has learned what the readings
hold.

\textit{The rate.} Gradient descent on the same fit at 0.9 of its
stability limit $2/\lambda_{\max}$ brings the training miss to
\SI{0.28}{\milli\electronvolt} in 400 steps; at 1.1 of it the miss grows
without bound, to \SI{1.7e32}{\milli\electronvolt}; at a thousandth of
it the miss has fallen only to \SI{3.6}{\milli\electronvolt}
(\cref{fig:lo-training}b). A loss that rises, or swings with growing
amplitude, means a rate past the limit; one that barely moves, a rate far
below it. Adam shows a milder form of the first: at a fixed rate of
$10^{-3}$ its training miss jumps by up to
\SI{0.24}{\milli\electronvolt} in a single step late in training, when
its running mean square has decayed and the steps no longer shrink with
the gradient, against \SI{0.005}{\milli\electronvolt} at $10^{-4}$, the
smaller rate removing the spikes.

\textit{Overfitting.} Degree 14 with no penalty, Adam at $10^{-4}$ for
\num{e5} steps: the training miss keeps falling, to
\SI{0.150}{\milli\electronvolt}, below the noise, while the validation
miss is least, \SI{0.38}{\milli\electronvolt}, after 160 steps and then
rises to \SI{2.95}{\milli\electronvolt} (\cref{fig:lo-training}c). Two
curves that part, the training miss below the noise and the validation
miss climbing, mean the model is learning the noise. Stopping at the
step of least validation miss, \emph{early stopping}, keeps the model of
that step and works as a penalty does, since the weights have not yet
grown into the directions the readings barely constrain; a penalty or
fewer weights cure it outright.

\textit{Unscaled features.} Degree 8 in $r$ itself, whose design has the
condition number \num{3.3e26} against \num{4.6e6} for the scaled
distance, trained by gradient descent at 0.9 of each one's stability
limit: after 2000 steps the training miss is
\SI{5.3}{\milli\electronvolt} with powers of $r$ and
\SI{0.22}{\milli\electronvolt} with the scaled distance
(\cref{fig:lo-training}d). A loss that falls a little and then crawls,
with no sign of the noise floor, points to badly scaled inputs before it
points to a bad model; scaling each input to a spread of about one is the
cure, and the reason the readings of a learned potential are normalised
before training.

\textit{A leaked split.} The fault that the loss curves cannot show is
the one of \cref{sec:lo-splits}: with frames split at random, training
and validation misses fall together to a healthy-looking floor, 30.5
against an independent run's \SI{128}{\milli\electronvolt}. Only a
validation set separated in time, or a separate run, reveals it.

\begin{figure}[htb]
\centering
\includegraphics{ch18_learning/training}
\caption{What the training looks like: the rms miss on the fifteen
training readings of argon's pair energy, and on fifteen validation
readings, against the step, from zero weights. (a) Healthy: degree 8,
Adam at $10^{-4}$. (b) Gradient descent at 0.9 and 1.1 of its stability
limit and at a thousandth of it. (c) Overfitting: degree 14, Adam at
$10^{-4}$, with the step of least validation miss (grey). (d) Degree 8 in
$r$ itself and in the scaled distance, gradient descent at 0.9 of each
one's limit. The noise of the readings, \SI{0.2}{\milli\electronvolt},
dotted.}
\label{fig:lo-training}
\end{figure}

\begin{table}[htb]
\centering
\small
\caption{Faults of training, the check that catches each, and its value
in the healthy run and in the broken one.}
\label{tab:lo-faults}
\begin{tabular}{@{}>{\raggedright\arraybackslash}p{0.22\linewidth}
   >{\raggedright\arraybackslash}p{0.34\linewidth}
   >{\raggedright\arraybackslash}p{0.17\linewidth}
   >{\raggedright\arraybackslash}p{0.17\linewidth}@{}}
\toprule
fault & check & healthy & broken\\
\midrule
rate past the limit & training miss after 400 steps & \SI{0.28}{\milli
\electronvolt} & \SI{1.7e32}{\milli\electronvolt}\\
rate far too small & training miss after 400 steps & \SI{0.28}{\milli
\electronvolt} & \SI{3.6}{\milli\electronvolt}\\
Adam's rate too large & largest jump of the training miss in a step &
\SI{0.005}{\milli\electronvolt} & \SI{0.24}{\milli\electronvolt}\\
overfitting & validation miss at the end against its least &
\num{0.19} against 0.18 (degree 8) & \num{2.95} against \SI{0.38}{\milli
\electronvolt}\\
unscaled features & condition number of $X^{\mathsf T}X$ & \num{4.6e6} &
\num{3.3e26}\\
leaked split & validation miss against an independent run's &
\num{92} against 128 (blocked) & \num{30.5} against
\SI{128}{\milli\electronvolt}\\
\bottomrule
\end{tabular}
\end{table}

\begin{keyidea}
The training and validation losses against the step are read as a run's
energies are: falling together to the noise is healthy; a rising or
swinging loss means a rate past the limit; a crawl, a rate too small or
inputs badly scaled; curves that part, overfitting, cured by early
stopping, a penalty or fewer weights. A leaked split looks healthy, and
only a validation set separated in time reveals it.
\end{keyidea}

\notebook{18}{break a training run in each of these ways and read its
loss curves}

\begin{selfcheck}
A model's training loss falls steadily, and its validation loss falls
for a while and then rises. What is happening, and which step's weights
should be kept?
\end{selfcheck}
